\section{Introducción}
\label{sec:intro}

Un cuadricóptero es un robot aéreo \emph{subactuado}: cuatro entradas (empuje total y tres torques) gobiernan seis grados de libertad, por lo que solo se pueden controlar cuatro salidas independientes. Las notas del curso resuelven este problema separando el movimiento en una dinámica de traslación, cuya entrada es el empuje $u_1$ proyectado por la actitud, y una dinámica de rotación, de la que se encarga un lazo interno rápido; ambas se linealizan por realimentación y se cierran en cascada~\cite{notas1}. La misma formulación advierte que la cancelación exacta depende de un modelo preciso y que las incertidumbres de $m$, $I$ o $k_f$, la saturación de los actuadores y las perturbaciones externas degradan el desempeño, lo que motiva el control adaptativo y robusto del segundo bloque del curso~\cite{notas2}.

Este informe presenta la implementación de esa teoría en un dron real: un marco F450 con cuatro motores A2212 de 1000\,KV, controlado por un ESP32-S3 con IMU MPU-6500, un sensor de distancia láser VL53L0X y un ultrasonido HC-SR04. El firmware evaluado es la versión V\_11 del proyecto. El objetivo es \emph{estabilizar la actitud (roll, pitch, yaw) y mantener la altura} usando un modelo cuyos parámetros se identificaron experimentalmente y evaluar el resultado con las cajas negras de vuelo. No se controla la posición lateral $(x,y)$, porque el dron no dispone de ningún sensor que la mida.

Las contribuciones son: (i) la identificación de masa, inercias y coeficientes del rotor (Sección~\ref{sec:ident}); (ii) un firmware de 250\,Hz que implementa la cascada del curso, un mezclador con prioridad de ejes y un estimador de altura por Kalman, con extensiones robusta y adaptativa (Sección~\ref{sec:control}); y (iii) el análisis cuantitativo de dos vuelos V\_11 y su comparación con V\_10, incluyendo los fallos observados (Sección~\ref{sec:res}). La Sección~\ref{sec:conc} resume las conclusiones.
